\FloatBarrier
\section{Appendix E. External Benchmark Details}
\phantomsection\label{app:e}
\subsection{Original-task Reruns and Scoring}
We reran BERT~\citep{devlin2019bert}, SpanBERT~\citep{joshi2020spanbert}, JobBERT, and JobSpanBERT~\citep{zhang2022skillspan} separately for skills and knowledge on SkillSpan's public HOUSE and TECH subsets. The MaChAmp~\citep{vandergoot2021machamp} \texttt{seq\_bio} CRF runs used 20 epochs, development-set checkpoint selection, and seeds 3477689, 4213916, 8749520, 6828303, and 9364029. BIG was unavailable, and we did not rerun the multi-task configuration. Table~\ref{tab:external-native} gives MaChAmp per-site span-F1. Table~\ref{tab:external-native-pooled} gives pooled HOUSE and TECH scores after BIO rescoring with the correct task column. These scoring procedures yield different results. The native runs contain 3,570 test sentences, compared with 3,569 in the common-protocol release; the one-record difference has not been resolved by matching IDs. For published test references, we use only the original paper's TEST/STL rows. Its HOUSE/TECH rows report development results.

For German public ICT, jobBERT-de used a 20-epoch MaChAmp configuration and ESCOXLM-R used its released five-epoch configuration~\citep{zhang2023escoxlmr}. Both used seeds 276800, 381552, 497646, 624189, and 884832. To run ESCOXLM-R with newer MaChAmp, we added the required \texttt{shuffle}, \texttt{diverse}, and \texttt{reset\_transformer\_model} fields without changing the learning rate or epoch limit. The runs therefore use the public releases with compatibility changes, rather than recreating every original environment. They contain 2,557 test sentences, whereas the ESCOXLM-R paper reports 2,943. They cover ICT extraction, not the complete German EDU/EXP/LNG pipeline~\citep{gnehm2022skills}.

\subsection{Data and Implementation Sources}
We obtained the common-protocol data from the public Hugging Face releases of
\href{https://huggingface.co/datasets/jjzha/skillspan}{SkillSpan},
\href{https://huggingface.co/datasets/jjzha/kompetencer}{Kompetencer},
\href{https://huggingface.co/datasets/jjzha/gnehm}{Gnehm ICT},
\href{https://huggingface.co/datasets/jjzha/green}{Green},
\href{https://huggingface.co/datasets/jjzha/sayfullina}{Sayfullina}, and
\href{https://huggingface.co/datasets/jjzha/fijo}{FIJO}.
The corresponding papers are cited in the evaluation protocol; preprocessing and splits in these releases may differ from the original studies. Code is available from the
\href{https://github.com/kris927b/SkillSpan}{SkillSpan},
\href{https://github.com/jjzha/kompetencer}{Kompetencer},
\href{https://github.com/mainlp/escoxlmr}{ESCOXLM-R}, and
\href{https://github.com/machamp-nlp/machamp}{MaChAmp} repositories.
The \href{https://github.com/jjzha/nnose}{NNOSE repository} supplies the retrieval-augmented implementation used in the subsequent A100 experiments.
Our linear-head protocol and compatibility changes are described above. Recorded revisions are listed in the experiment records below.

\subsection{Uncertainty and Experiment Records}
Table~\ref{tab:external-common-ci} gives the paired sentence-bootstrap intervals. Sayfullina's interval includes zero, although its lower endpoint is close to zero. Because the intervals are pointwise, exclusions of zero do not constitute significance tests adjusted for multiple comparisons. The intervals capture sentence sampling uncertainty conditional on the trained seeds; seed SD describes variation across runs. Neither accounts for dependence between sentences from the same advertisement.

The \href{https://github.com/AlfredJamesLi/chinese-skillspan-benchmark/tree/46d1166f8e11dd52d473f645536074bb8dd0a53c/results_snapshots/external_benchmarks_20260922}{server-B experiment records} contain 104 completed training runs: 48 under the common protocol, six for Chinese transfer, 40 for SkillSpan, and ten for German ICT. We used XLM-R-large and ESCOXLM-R at repository revisions identified by the abbreviated commit hashes \texttt{c23d21b0} and \texttt{8093cc37}, respectively. The archive records the base-model revisions, but some dataset and executed-code identifiers remain unresolved. We recomputed means and sample SD from per-seed metrics and checked the paired differences and interval records for consistency. Table summaries use the unrounded per-seed metrics. These checks used the reported outputs; they did not independently rerun training, rescore predictions, or regenerate intervals. The common-protocol and Chinese records list prediction and configuration paths and hashes, but those entries alone do not provide the underlying files for independent verification. The linked audit documents missing execution and environment records; the archive is not a complete training package.

All seeds remain in the encoder summaries. ESCOXLM-R selected epoch 1 for SkillSpan knowledge at seed 43 and epoch 2 for German ICT at seed 42. The largest mean decreases in the common-protocol comparison were 0.027 on SkillSpan knowledge and 0.023 on German ICT.

The subsequent \href{https://github.com/AlfredJamesLi/chinese-skillspan-benchmark/tree/b84f218b7142c97ab03a60499a9b30e7ec332cab/results_snapshots/a_native_20260924}{server-A archive} adds 30 Kompetencer classification training runs and six NNOSE training runs. Its 72 metric rows include two evaluation outputs per trained model; they are not 72 separate training runs. Classification means and sample SD were recomputed from the five-seed records, and NNOSE values and prediction hashes were checked for agreement between the summary and the result ledger. These checks establish consistency among the reported summaries; they do not independently rescore the predictions. The archive lists model revisions, software versions, and per-seed results.

\begin{table}[H]
\centering\small
\caption{Paired differences under the common protocol (ESCOXLM-R minus XLM-R).}
\label{tab:external-common-ci}
\begin{tabularx}{\linewidth}{@{}L P{.19\linewidth}rr@{}}
\toprule
Dataset & Stream & Mean difference & 95\% interval \\
\midrule
SkillSpan & Skill & $+0.010$ & $[-0.003, 0.023]$ \\
SkillSpan & Knowledge & $-0.027$ & $[-0.041, -0.013]$ \\
Kompetencer & Skill & $-0.018$ & $[-0.063, 0.022]$ \\
Kompetencer & Knowledge & $-0.025$ & $[-0.102, 0.056]$ \\
Gnehm ICT & ICT & $-0.023$ & $[-0.034, -0.013]$ \\
Green & Native types & $-0.001$ & $[-0.021, 0.017]$ \\
Sayfullina & Skill & $+0.006$ & $[-1.47\times10^{-5}, 0.011]$ \\
FIJO & Native types & $-0.004$ & $[-0.043, 0.038]$ \\
\bottomrule
\end{tabularx}
\par\smallskip\RaggedRight\footnotesize Pointwise sentence-bootstrap intervals (2,000 resamples; seed 20260921), conditional on seeds 42/43/44. No advertisement clustering or multiple-comparison adjustment. Values are rounded after computing differences and intervals. Three decimal places are used except for Sayfullina's near-zero lower endpoint, shown with three significant digits in scientific notation to preserve its negative sign. Its interval includes zero.
\end{table}

\begin{table}[H]
\centering\small
\caption{SkillSpan: published single-task TEST references and public-subset MaChAmp reruns (0--1; five seeds).}
\label{tab:external-native}
\begin{tabularx}{\linewidth}{@{}L P{.15\linewidth}rrr@{}}
\toprule
Model & Task & Published TEST & Rerun HOUSE & Rerun TECH \\
\midrule
BERT & Skills & $0.543\pm0.007$ & $0.496\pm0.013$ & $0.514\pm0.012$ \\
BERT & Knowledge & $0.624\pm0.004$ & $0.550\pm0.004$ & $0.700\pm0.008$ \\
SpanBERT & Skills & --- & $0.516\pm0.014$ & $0.526\pm0.014$ \\
SpanBERT & Knowledge & --- & $0.536\pm0.017$ & $0.691\pm0.015$ \\
JobBERT & Skills & $0.561\pm0.005$ & $0.534\pm0.014$ & $0.538\pm0.013$ \\
JobBERT & Knowledge & $0.639\pm0.003$ & $0.577\pm0.004$ & $0.708\pm0.005$ \\
JobSpanBERT & Skills & $0.566\pm0.008$ & $0.543\pm0.013$ & $0.519\pm0.010$ \\
JobSpanBERT & Knowledge & $0.611\pm0.010$ & $0.546\pm0.018$ & $0.687\pm0.014$ \\
\bottomrule
\end{tabularx}
\par\smallskip\RaggedRight\footnotesize Published values are the TEST/STL results in Table 5 of \citet{zhang2022skillspan}, converted from percent and rounded to three decimal places; their HOUSE/TECH rows report development results. Our reruns use MaChAmp per-site span-F1 on public HOUSE+TECH, excluding BIG. The evaluation samples differ, so score differences cannot be interpreted as matched replication gaps. The source table gives no SpanBERT TEST score. Multi-task training was not rerun.
\end{table}

\begin{table}[H]
\centering\small
\caption{SkillSpan pooled HOUSE+TECH BIO rescoring of the same native runs.}
\label{tab:external-native-pooled}
\begin{tabularx}{\linewidth}{@{}Lrr@{}}
\toprule
Model & Skills & Knowledge \\
\midrule
BERT & $0.495\pm0.011$ & $0.639\pm0.005$ \\
SpanBERT & $0.511\pm0.012$ & $0.632\pm0.009$ \\
JobBERT & $0.525\pm0.008$ & $0.654\pm0.003$ \\
JobSpanBERT & $0.522\pm0.008$ & $0.628\pm0.008$ \\
\bottomrule
\end{tabularx}
\par\smallskip\RaggedRight\footnotesize Mean $\pm$ sample SD across five seeds, on a 0--1 scale. Predictions were rescored using the correct BIO task column and pooled across HOUSE+TECH. These scores differ from averaging site F1 or applying the per-site MaChAmp scorer above. The native runs contain 3,570 sentences, compared with 3,569 in the common-protocol release; the one-record difference remains unresolved at the ID level.
\end{table}

\begin{table}[H]
\centering\small
\caption{German public ICT extraction: native-task recipes (0--1; five seeds).}
\label{tab:external-gnehm}
\begin{tabularx}{\linewidth}{@{}L P{.25\linewidth}rr@{}}
\toprule
Task/model & Recipe & Published & Rerun \\
\midrule
Public ICT: jobBERT-de & 20-epoch MaChAmp & --- & $0.885\pm0.010$ \\
Public ICT: ESCOXLM-R & 5-epoch released recipe & $0.884\pm0.005$ & $0.891\pm0.005$ \\
\bottomrule
\end{tabularx}
\par\smallskip\RaggedRight\footnotesize Reruns use 2,557 test sentences. The published ESCOXLM-R score comes from Table 2 of \citet{zhang2023escoxlmr}, whose Table 1 reports 2,943 test sentences. These different releases prevent a matched comparison with the published score. No corresponding published jobBERT-de score is supplied. Epoch budgets differ between the two models here, unlike in the common-protocol comparison.
\end{table}


\FloatBarrier
\subsection{Kompetencer Classification and NNOSE Reruns}
Kompetencer fine-grained classification~\citep{jensen2022kompetencer} and NNOSE~\citep{zhang2024nnose} were subsequently run on an NVIDIA A100 (80 GB) using CUDA 11 software builds. The older stacks had not run on server B's Blackwell GPU. These experiments retain their task-specific metrics and are separate from the common linear-head protocol.

For Kompetencer, we used MaChAmp with AllenNLP 2.8.0 and PyTorch 1.8.1+cu111, a 20-epoch budget, and the same five seeds listed for the SkillSpan native runs. The training conditions used Danish (DA), English (EN), or their combination (EN+DA); the source files contain 138 Danish and 9,472 English training records. Development selection used the 1,577-record English development set. Table~\ref{tab:external-kompetencer-classification} reports scikit-learn macro-F1 on 784 Danish and 1,578 English classified spans, with sample SD across five seeds. This evaluates classification of supplied spans, rather than their extraction. The requested DaJobBERT cased model identifier redirected to an uncased checkpoint; consequently, these runs do not recreate a cased configuration. The source plotting code labels its metric ``Weighted macro-F1'', which differs from the unweighted macro-F1 reported here. Weighted scores remain in the experiment archive, and a matched comparison with the original published matrix has not been established.

For NNOSE, we used PyTorch 1.10.1+cu113, adapter-transformers 3.2.0, and FAISS 1.7.2 with JobBERT and JobBERTa. Each model used seed 113412 and development-based early stopping. The datastore contained only training examples from the dataset being evaluated. Neighbour count $k$, interpolation weight $\lambda$, and temperature $T$ were selected on development data and frozen before test evaluation. Table~\ref{tab:external-nnose} reports seqeval span-F1 for both the trained linear head and its retrieval-augmented predictions. These linear-head scores come from the NNOSE pipeline, not the common-protocol XLM-R experiments. The NNOSE release contains 3,569 SkillSpan, 335 Green, and 1,851 Sayfullina test sentences. The all-dataset datastore configuration ($\forall D$) was not run. All six conditions use one seed, so no seed SD or general retrieval advantage is inferred.
\begin{table}[htbp]
\centering\small
\caption{Kompetencer fine-grained classification: macro-F1 across five seeds (0--1).}
\label{tab:external-kompetencer-classification}
\begin{tabularx}{\linewidth}{@{}L P{.18\linewidth}rr@{}}
\toprule
Encoder & Training & Danish test & English test \\
\midrule
DaJobBERT & DA & $0.182\pm0.034$ & $0.064\pm0.012$ \\
DaJobBERT & EN & $0.232\pm0.028$ & $0.539\pm0.010$ \\
DaJobBERT & EN+DA & $0.340\pm0.017$ & $0.529\pm0.008$ \\
RemBERT & DA & $0.050\pm0.072$ & $0.034\pm0.041$ \\
RemBERT & EN & $0.261\pm0.019$ & $0.552\pm0.017$ \\
RemBERT & EN+DA & $0.316\pm0.019$ & $0.554\pm0.014$ \\
\bottomrule
\end{tabularx}
\par\smallskip\RaggedRight\footnotesize Mean $\pm$ sample SD; Danish/English test sizes are 784/1,578 classified spans. All conditions use English development data. EN training with Danish testing is the EN-to-DA transfer condition. DaJobBERT resolved to an uncased checkpoint despite the requested cased identifier. Scores are unweighted classification macro-F1, not span-F1 or the source plot's weighted metric. No matched published-score comparison is claimed.
\end{table}

\begin{table}[htbp]
\centering\small
\caption{NNOSE with an in-dataset training datastore: test span-F1 (0--1).}
\label{tab:external-nnose}
\begin{tabularx}{\linewidth}{@{}L P{.15\linewidth}r r r r P{.19\linewidth}@{}}
\toprule
Dataset & Encoder & Test $n$ & Epoch & Linear & kNN & $(k,\lambda,T)$ \\
\midrule
SkillSpan & JobBERT & 3,569 & 17 & 0.613 & 0.612 & $(16,0.25,3)$ \\
SkillSpan & JobBERTa & 3,569 & 6 & 0.631 & 0.625 & $(8,0.4,0.5)$ \\
Green & JobBERT & 335 & 1 & 0.491 & 0.500 & $(128,0.2,2)$ \\
Green & JobBERTa & 335 & 3 & 0.466 & 0.485 & $(128,0.5,0.5)$ \\
Sayfullina & JobBERT & 1,851 & 7 & 0.897 & 0.896 & $(16,0.85,0.5)$ \\
Sayfullina & JobBERTa & 1,851 & 6 & 0.915 & 0.912 & $(128,0.6,0.1)$ \\
\bottomrule
\end{tabularx}
\par\smallskip\RaggedRight\footnotesize Single seed 113412; no SD is estimated. Epoch identifies the selected checkpoint; $k$, $\lambda$, and $T$ are development-selected retrieval settings. Both columns use the same trained checkpoint within each row. The datastore excludes development and test examples. SkillSpan uses public HOUSE+TECH. The all-dataset datastore was not evaluated, and alignment with an original-paper result cell has not been established.
\end{table}

In the classification runs, combined English and Danish training yielded higher Danish macro-F1 than Danish-only training for both encoders. RemBERT's Danish-only scores varied substantially across seeds ($0.050\pm0.072$). These comparisons describe the recorded training conditions and English-based model selection, not a matched replication of the published weighted metric. In NNOSE, retrieval increased test F1 on Green for both encoders but slightly decreased it on SkillSpan and Sayfullina. This single-seed result does not establish a consistent benefit from retrieval.

BERT, JobBERT, and DaBERT classification configurations, NNOSE's all-dataset datastore, the full German EDU/EXP/LNG pipeline, commercial-API comparisons, and SkillSpan multi-task training remain without completed results here. They are not assigned zero scores. These external experiments broaden the method coverage but do not replace an independent Chinese human test.

\FloatBarrier
